\documentclass[11pt,a4paper]{article}

\usepackage[T1]{fontenc}
\usepackage[utf8]{inputenc}
\usepackage[a4paper,top=2.4cm,bottom=2.4cm,left=2.6cm,right=2.6cm]{geometry}
\usepackage{amsmath,amssymb}
\usepackage{booktabs,array,tabularx}
\usepackage{enumitem}
\usepackage{xcolor}
\usepackage{listings}
\usepackage{titlesec}
\usepackage{fancyhdr}
\usepackage{parskip}

\definecolor{ink}{RGB}{28,32,38}
\definecolor{accent}{RGB}{20,72,168}
\definecolor{codebg}{RGB}{244,246,249}
\definecolor{conbg}{RGB}{29,33,40}
\definecolor{congreen}{RGB}{150,230,170}
\definecolor{rule}{RGB}{210,216,224}
\definecolor{cmt}{RGB}{110,120,132}

\usepackage[colorlinks=true,linkcolor=accent,urlcolor=accent,citecolor=accent]{hyperref}

\titleformat{\section}{\Large\bfseries\color{ink}}{\thesection.}{0.5em}{}
\titleformat{\subsection}{\large\bfseries\color{ink}}{\thesubsection}{0.5em}{}
\titlespacing*{\section}{0pt}{1.4em}{0.5em}
\titlespacing*{\subsection}{0pt}{1.0em}{0.3em}

\pagestyle{fancy}\fancyhf{}
\renewcommand{\headrulewidth}{0.4pt}
\renewcommand{\headrule}{\hbox to\headwidth{\color{rule}\leaders\hrule height \headrulewidth\hfill}}
\fancyhead[L]{\small\color{cmt}GARF for Stata — Installation \& Use Guide}
\fancyhead[R]{\small\color{cmt}\thepage}

% ---- code (Stata) ----
\lstdefinestyle{stata}{
  backgroundcolor=\color{codebg}, basicstyle=\ttfamily\small\color{ink},
  breaklines=true, breakatwhitespace=true, columns=fullflexible,
  keepspaces=true, frame=single, framerule=0pt, rulecolor=\color{codebg},
  xleftmargin=8pt, xrightmargin=8pt, aboveskip=6pt, belowskip=6pt,
  commentstyle=\color{cmt}\itshape, morecomment=[l]{*}, morecomment=[l]{//},
}
% ---- console (dark) ----
\lstdefinestyle{console}{
  backgroundcolor=\color{conbg}, basicstyle=\ttfamily\small\color{congreen},
  breaklines=true, columns=fullflexible, keepspaces=true,
  frame=single, framerule=0pt, rulecolor=\color{conbg},
  xleftmargin=8pt, xrightmargin=8pt, aboveskip=6pt, belowskip=6pt,
}
\lstnewenvironment{stata}{\lstset{style=stata}}{}
\lstnewenvironment{console}{\lstset{style=console}}{}
\newcommand{\opt}[1]{\texttt{\textbf{#1}}}
\newcommand{\code}[1]{\texttt{#1}}

\begin{document}

% ---------------- title ----------------
\thispagestyle{empty}
\begin{center}
{\Huge\bfseries\color{ink} GARF for Stata}\\[0.4em]
{\Large\color{accent} Installation \& Use Guide}\\[1.0em]
{\large The \texttt{garf} and \texttt{garfte} commands — forest-proximity matching for panel counterfactuals, and the effect estimates that follow}\\[0.4em]
{\color{cmt}\small Binary build · \textbf{version 1.2} · this guide covers \texttt{garf.ado} and \texttt{garfte.ado}}\\[0.7em]
{\small Federico Fantechi, PhD · Department of Economics, Business and Statistics · University of Palermo}\\[0.15em]
{\small \texttt{federico.fantechi@unipa.it}}
\end{center}
\vspace{0.6em}
\hrule height 0.4pt
\vspace{0.4em}

\begin{center}
\fcolorbox{accent}{codebg}{\parbox{0.94\textwidth}{\small
\textbf{New in version 1.2: the estimation stage.} The companion command \code{garfte} runs
after \code{garf} and estimates the effect on the treated from the matched panel: a matching
difference-in-differences (\opt{method(did)}) or the treated-minus-own-donors gap in levels
(\opt{method(gap)}), by event time and overall, with an event-study graph, \code{e(b)}/\code{e(V)},
and the effect of each treated unit in a file (\S\ref{sec:garfte}). \code{garf} computes exactly as
in 1.1; it now stores its settings on the loaded panel (characteristics \code{\_dta[garf\_*]}, read
by \code{garfte}), imports the matched panel in double precision, and ends with a short ``Next
steps'' note.

\medskip
\textbf{Version 1.1} is the first public release. It is the development build beta~1.5 with three changes to the Stata command: names and paths are quoted in the configuration file passed to the engine, the data are exported without value labels, and the temporary copy of the data is erased after the run. The notes below keep the development numbering.

\medskip
\textbf{New in beta 1.5.} Forest Match Quality is normalised by the best donor \emph{attainable}
in the data, so that 1 means ``every donor as close as the closest untreated unit available''
and its level no longer depends on the sample size or the leaf size (\S\ref{sec:fmq}); the
earlier version is still returned as \code{r(fmq\_twin\_mean)}. In \code{DIAGNOSTICS}, the widest
design of the sensitivity grid is the setting \code{KWIDE}, default \code{k(1000)}.

\smallskip
\textbf{New in beta 1.4.} No change to the command or to the computations: results are
identical to beta 1.3. The package now ships the folder \code{DIAGNOSTICS}: commented do-files,
one per diagnostic, for the full design-stage protocol (\S\ref{sec:diagnostics}). It replaces the
old testing do-files.

\smallskip
\textbf{New in beta 1.3.} A leaner encoding: each time-varying variable now gives $L+7$
features instead of $L+9$, because the first and last level of the window, exact copies of two
lags, are no longer repeated (results therefore differ slightly from beta 1.2). Growth is now
\emph{symmetric} growth, (last $-$ first)$/$((|last| + |first|)$/2$), defined for zeros and negative
values. The forest's fit is now reported honestly, from its \emph{out-of-bag} predictions:
balanced accuracy \code{r(oob\_bacc)} and AUC \code{r(oob\_auc)}; the in-sample \code{r(auc)} was
removed. The legacy option \opt{minprox()} (use \opt{caliper()}) and the unused option
\opt{load} were removed.

\smallskip
\textbf{New in beta 1.2.} \opt{wpower(\#)} sets the donor weights to
$\text{proximity}^{\#}$, normalised within each treated unit (\S\ref{sec:proxcalc}); the default
\opt{wpower(1)} reproduces the beta 1.1 weights exactly, \opt{wpower(0)} gives uniform $k$:1
weights. Engine fixes: \opt{noreplacement} now matches the hardest-to-match treated units first;
the matched $|$SMD$|$ now compares only the matched treated units with their donors. The
time-invariant option now abbreviates to \opt{tin()} (it clashed with \opt{time()}). New
diagnostic: \textbf{Forest Match Quality (FMQ)}, \S\ref{sec:fmq} --- it systematically rewards fewer,
closer donors, so read it together with ESS.

\smallskip
\textbf{Changed in beta 1.1.} \opt{caliper()} now filters on a
\textbf{min-max normalised} proximity, on a fixed $[0,1]$ scale, instead of the raw
leaf-proximity kernel used in beta 1.0 --- see \S\ref{sec:proxcalc}. \textbf{A caliper value
chosen under beta 1.0 does not carry over}; re-pick it using the new \code{proximity\_minmax}
column/variable. The raw \code{proximity} column is unchanged.}}
\end{center}
\vspace{0.4em}

{\small
\textbf{In one line.} \code{garf} runs the \emph{Group Aware Random Forest} Python engine from
inside Stata and brings the results back: you give it a panel and it returns a \emph{matched
donor panel} ready for a weighted difference-in-differences; \code{garfte} then estimates the
effect on that panel.

\smallskip
\noindent The \emph{Group Aware Random Forest} (GARF) identifies matching (untreated)
observations over a very large set of covariates, where linear models (e.g.\ propensity scores)
often fail; \emph{group awareness} lets the matching run over a panel structure --- matching not
only the \emph{levels} of each covariate but also their \emph{dynamics}; and GARF is highly
flexible, adjustable for issues such as a low number of treated units and feature selection, with
parameters that can be tuned for better precision.

\smallskip
\noindent \code{garf} itself does not estimate the treatment effect: the companion command
\code{garfte} does (\S\ref{sec:garfte}), or you take the matched panel to your own estimator.}

\vspace{0.4em}
\hrule height 0.4pt
\vspace{0.8em}

{\small\tableofcontents}
\vspace{0.6em}

% =====================================================================
\section{What this is, and how it works}

\code{garf} is a thin Stata wrapper (\code{garf.ado}) around the Python engine
\code{garf\_binary.py} that ships with the \textbf{GARF-binary} package. Stata does not
re-implement the method; it drives the same Python code, so \emph{no features are lost} ---
including optional hyperparameter tuning.

When you run \code{garf}, the command:
\begin{enumerate}[leftmargin=1.5em,itemsep=1pt,topsep=2pt]
\item exports the panel currently in memory (honouring \code{if}/\code{in}) to a temporary CSV;
\item writes a configuration file from the options you passed;
\item shells out to the Python engine, which fits the random forest, builds the leaf-proximity
      kernel, and selects each treated unit's \code{k} nearest control ``donors'';
\item reads the results back and prints a matching summary, storing key numbers in \code{r()};
\item \textbf{by default loads the matched panel into memory} (saving no files) so you can run
      the DiD immediately; \code{store} / \code{export()} additionally write the CSV outputs.
\end{enumerate}

The division of labour is deliberate: \textbf{\code{garf} selects the comparison group; the
effect is estimated in a separate step}, either with the package's \code{garfte} (a second Python
file, \code{garf\_estimate.py}, driven the same way; \S\ref{sec:garfte}) or with your own weighted
DiD (\code{reghdfe}, \code{xtreg}, \code{didregress}). The design can thus be fixed, and checked
with the diagnostics, before any effect is looked at.

\medskip
\noindent\textbf{Why the shell-out design.} It reuses the package's own Python environment
(\code{.venv}), so you do \emph{not} have to configure Stata's built-in Python. Any Stata that
can run \code{shell} works, on Windows, macOS, or Linux.

% =====================================================================
\section{Requirements}

\begin{itemize}[leftmargin=1.5em,itemsep=2pt,topsep=2pt]
\item \textbf{Stata 14 or newer} (uses \code{export/import delimited}, \code{shell},
      \code{file}, \code{regexm}).
\item \textbf{Python 3.10+} with the package dependencies. The easiest way to get them is to
      run the package's \code{setup} once (next section): it builds a local \code{.venv} that
      \code{garf} finds automatically.
\item The \textbf{GARF-binary package} folder, unzipped somewhere stable.
\item For the DiD step, \code{reghdfe} is convenient (\code{ssc install reghdfe}); plain
      \code{xtreg}/\code{areg} also work.
\end{itemize}

You do \emph{not} need Stata's Python integration (\code{python:}); \code{garf} uses
\code{shell}, not \code{python:}.

% =====================================================================
\section{Installation (one-time)}

\subsection{Step 1 --- set up the Python environment}
From the \textbf{GARF-binary} folder, run the setup once. This creates \code{.venv} and
installs the dependencies; nothing is installed system-wide.
\begin{stata}
* macOS / Linux (Terminal, inside the GARF-binary folder):
./setup.sh

* Windows: double-click  setup.bat
\end{stata}
If you do not have Python yet, follow \code{INSTALL.txt} in the package first.

\subsection{Step 2 --- put garf.ado on Stata's ado-path}
Two choices.

\textbf{(a) Just for this session} --- add the folder to the ado-path:
\begin{stata}
adopath + "C:/Users/you/GARF-binary/integrations/stata"   // Windows
adopath + "/Users/you/GARF-binary/integrations/stata"     // macOS / Linux
\end{stata}

\medskip
\noindent{\color{accent}\textbf{The path must end in \code{/stata}}} --- the folder that
actually contains \code{garf.ado} and \code{garf.sthlp}. Stata's ado-path is \emph{not}
recursive: adding the parent \code{.../integrations} will \emph{not} work, and both
\code{help garf} and the \code{garf} command will appear to be missing. Confirm it is found
with:
\begin{stata}
which garf     // should print the path to garf.ado
\end{stata}

\textbf{(b) Permanently} --- copy \code{garf.ado}, \code{garf.sthlp}, \code{garfte.ado} and
\code{garfte.sthlp} into your PERSONAL ado folder. Find it with:
\begin{stata}
sysdir
\end{stata}
and copy the four files into the folder shown next to \code{PERSONAL} (create it if needed).
Stata will then find \code{garf} in every session. Keep the \code{GARF-binary} folder itself in a
fixed place: \code{garf} finds the engine through \code{GARF\_HOME} (Step 3).

\medskip
\noindent{\color{accent}\textbf{Check which version runs.}} Stata looks for \code{garf.ado} in
the \emph{current} folder before the ado-path, so an old copy left in a working folder silently
takes precedence over the one in PERSONAL. After installing or updating, run
\begin{stata}
discard        // drop any garf already loaded in this session
which garf     // the first line must read: version 1.2
which garfte   // the first line must read: version 1.2
\end{stata}

\subsection{Step 3 --- point Stata at the package}
\code{garf} needs to know where the engine lives. Set a global once:
\begin{stata}
global GARF_HOME "C:/Users/you/GARF-binary"   // Windows (forward slashes are fine)
global GARF_HOME "/Users/you/GARF-binary"     // macOS / Linux
\end{stata}
To make it automatic, add that \code{global} line (and, if you chose 2(a), the \code{adopath}
line) to your \code{profile.do}, which Stata runs at startup.

\subsection{Step 4 --- verify}
\begin{stata}
which garf     // prints the path to garf.ado if the ado-path is right
help garf      // opens the help file
help garfte    // the estimation stage
\end{stata}
If either comes up empty, the ado-path line is pointing at the wrong folder --- re-check that it
ends in \code{/stata} (previous step). For an end-to-end check, run the package demo once (it writes
\code{demo/input/panel.csv}), then:
\begin{stata}
import delimited "$GARF_HOME/demo/input/panel.csv", clear varnames(1) case(preserve)
garf, treatment(W) outcome(Y) id(id) time(year) t0(2000) tvarying(x1 x2) tinvariant(region)
\end{stata}

% =====================================================================
\section{Quick start}

With a long panel in memory (one row per unit-year, a 0/1 treatment column that is 1 in the
treated unit's treated years):
\begin{stata}
garf, treatment(W) outcome(Y) id(id) time(year) t0(2000) ///
      tvarying(x1 x2) tinvariant(region) k(50)
\end{stata}
You will see a summary like this (numbers are from the synthetic demo panel):
\begin{console}
[garf] engine : /Users/you/GARF-binary/engine/garf_binary.py
[garf] python : /Users/you/GARF-binary/.venv/bin/python
[garf] tuning : off   k=50   caliper=0 (normalised)   wpower=1   replacement=true
[garf] running... (this can take a while when tuning is ON)
--------------------------------------------------------------
GARF (binary) - matching summary
--------------------------------------------------------------
Treated units              =        60
Control units              =       240
Treated matched            =        60
Distinct donors used       =       129
Mean |SMD| (raw -> matched)= 1.013  ->  0.428
Max  |SMD| (raw -> matched)= 2.197  ->  0.893
Mean ESS per treated       =   5.6 / 50
Match quality FMQ mean/p10 = 0.630 / 0.556   (all treated: 0.630)
  FMQ: 0 = donors no closer than random controls, 1 = as close as the best attainable donor.
  Fewer/closer donors systematically raise it: compare designs together with ESS.
Forest OOB bal.acc / AUC   = 0.958 / 0.993   (OOB accuracy 0.963)
  Near 0.5 = treatment hard to predict (good overlap); near 1 = groups separable.
  With rare treatment read the AUC first: balanced accuracy can stay near 0.5.
--------------------------------------------------------------
outputs: /Users/you/current-folder/garf_output
--------------------------------------------------------------
\end{console}
Read the balance line as the headline: matching cut the mean absolute standardized
difference from 1.10 to 0.42. The effective sample size (ESS) per treated unit is well below
\code{k} because good donors are reused --- that is expected, and worth reporting.

By default the matched panel is now \emph{in memory} and nothing was written to disk; the next
sections show the variables it adds and how to estimate the DiD. This replaces the data in
memory, so \code{preserve} or save first; add \code{store} to also keep the CSVs.

The output ends with a short note on what to do next:
\begin{console}
--------------------------------------------------------------
Next steps
  Save the matched panel (it replaced your data; garfte's settings are saved with it):
     save "garf_matched.dta", replace
  Diagnostics (design stage, pre-treatment only): see help garf, Diagnostics
     with this matched panel in memory (its settings are read from it, nothing to edit):
     do "/Users/you/GARF-binary/DIAGNOSTICS/00_master.do"   (balance, placebo, overlap, FMQ, sensitivity)
  Estimation: garfte, after garf
     garfte, method(gap)   treated minus its own donors, in levels
     method(did), the matching DiD, needs a panel matched without replacement:
     rerun garf with noreplacement (and a k the donor pool can supply), then garfte, method(did)
--------------------------------------------------------------
\end{console}
Then estimate the effect (\S\ref{sec:garfte}):
\begin{stata}
garfte, method(gap)
\end{stata}

% =====================================================================
\section{Command reference}

\subsection{Syntax}
\begin{stata}
garf [if] [in], treatment(varname) outcome(varname) id(varname)
     time(varname) t0(#) [ options ]
\end{stata}

\subsection{Required options}
\begin{tabularx}{\textwidth}{@{}lX@{}}
\toprule
Option & Meaning\\
\midrule
\opt{treatment(varname)} & binary treatment column (1 in treated years, else 0)\\
\opt{outcome(varname)}   & the outcome variable. For the matching it is simply the first time-varying variable, treated exactly like those in \opt{tvarying()} (same pre-window features; the forest is fit on the treatment, never on the outcome). Naming it serves reporting and later steps: \code{garfte} estimates on it by default, the diagnostics plot its path and placebo gap. No need to repeat it in \opt{tvarying()}\\
\opt{id(varname)}        & unit identifier\\
\opt{time(varname)}      & time variable\\
\opt{t0(\#)}             & treatment onset period $T_0$\\
\bottomrule
\end{tabularx}

\subsection{Design options}
\begin{tabularx}{\textwidth}{@{}lX@{}}
\toprule
Option & Meaning (default)\\
\midrule
\opt{tvarying(varlist)}   & covariates that change over time; get lags \emph{and} trajectory features\\
\opt{tinvariant(varlist)} & covariates that do not change (strings/categoricals allowed --- no encoding)\\
\opt{npre(\#)}            & pre-window length used for matching: $[T_0-\text{npre},\,T_0-1]$ (5)\\
\opt{includet0}          & include the onset year in the pre-window (off)\\
\bottomrule
\end{tabularx}

\subsection{Matching options (donor selection)}
Two rules pick each treated unit's controls, and they \textbf{combine} (see the note below).

\begin{tabularx}{\textwidth}{@{}lX@{}}
\toprule
Option & Meaning (default)\\
\midrule
\opt{k(\#)}          & keep at most this many nearest donors --- k-NN (50)\\
\opt{caliper(\#)}    & keep only donors with \emph{normalised} proximity $\ge$ \# (range $[0,1]$;
\code{0} = off). Composes with \code{k}\\
\opt{nokcap}         & lift the \code{k} cap: keep \emph{all} donors within the caliper\\
\opt{noreplacement}  & match without replacement, hardest-to-match treated units first (default is \emph{with} replacement)\\
\opt{wpower(\#)}     & donor weights $\propto \text{proximity}^{\#}$, summing to 1 per treated unit (1; 0 = uniform $k$:1; $>1$ = more weight on the closest donors)\\
\bottomrule
\end{tabularx}

\medskip
\noindent\textbf{k-NN vs.\ caliper.} By default GARF keeps each treated unit's \code{k} nearest
donors. A \emph{caliper} instead keeps every donor whose \emph{min-max normalised} proximity
(\S\ref{sec:proxcalc}) is at least \code{caliper}, so the number of donors varies across treated
units --- and can fall to zero for a unit with no good match, which is an honest overlap signal,
not a bug. A donor must also have actually co-occurred with the treated unit at least once (raw
proximity $>0$); that floor applies regardless of the caliper. The two rules compose:
\code{caliper(0.3) nokcap} keeps all donors within the caliper; \code{caliper(0.3) k(50)} keeps
the 50 nearest \emph{among} those within it. A tighter (larger) caliper trades donors for quality
--- fewer, closer controls, better balance, smaller effective samples --- so always report how
many survive (the ESS line, \S\ref{sec:interpret}). Because the caliper is on a fixed $[0,1]$
scale rather than a forest-size-dependent raw number, \code{caliper(0.3)} means the same thing
--- a pair at least 30\% of the way from the least similar to the most similar pair observed in
this run --- across different specifications; see \S\ref{sec:proxcalc} for exactly how that
scale is built. (This is a rescaling of the \emph{range}, not a percentile: because most pairs
tend to sit near the low end and a few near the top, do not read \code{caliper(0.3)} as
``keeps 70\% of donors'' --- always check how many actually survive.)

\subsection{How proximity, weight, and the normalised proximity are calculated}\label{sec:proxcalc}

Three related numbers come out of the matching step, and they answer different questions.
Getting them mixed up is the easiest way to misread a GARF run, so this section gives all three
formulas explicitly.

\medskip
\noindent\textbf{1. Proximity (raw) --- ``how close is this pair, as the forest sees it.''}
For treated unit $i$ and control $j$, over the $B$ trees of the fitted forest,
\[
\text{proximity}(i,j) \;=\; \frac{1}{B}\sum_{b=1}^{B}
\frac{\mathbf{1}\{\text{leaf}_b(i) = \text{leaf}_b(j)\}}{\#\{\text{controls in that leaf of tree } b\}}.
\]
Push $i$ and $j$ down every tree; each tree contributes $1/(\text{leaf size})$ if they land in the
same terminal leaf, $0$ otherwise; average over trees. This is Breiman's leaf-proximity idea, read
as an adaptive nearest-neighbour kernel (Lin and Jeon, 2006): a pair that keeps landing together
in small, tight leaves is close under the metric the forest actually learned. It is $0$ for a pair
that never once shares a leaf, and \textbf{higher means closer} --- the opposite convention from a
distance. Its absolute size is \emph{not} comparable across runs: it shrinks as \opt{minleaf()}
falls, moves with \opt{nestimators()}, and depends on how many covariates are in play, so
\code{proximity = 0.03} can be a strong match in one specification and a weak one in another.

\medskip
\noindent\textbf{2. Proximity, min-max normalised --- ``how close is this pair, on a fixed
0--1 scale.''} New in beta 1.1. The raw proximity matrix for the whole run is rescaled
\emph{once}, \emph{globally}:
\[
\text{proximity\_minmax}(i,j) \;=\; \frac{\text{proximity}(i,j) - \min}{\max - \min},
\]
where $\min$ and $\max$ are the smallest and largest raw proximity value anywhere in the entire
treated~$\times$~control matrix produced by that run (reported as \code{proximity\_raw\_min} /
\code{proximity\_raw\_max} in \code{diagnostics.json}). This is a \emph{global} rescaling, not a
per-treated-unit one, and that is a deliberate choice, not an implementation shortcut:
normalising each treated unit's row by its \emph{own} maximum would send every unit's
best-available donor to $1.0$ regardless of how poor that donor actually is, which would erase
the one signal that tells you a treated unit has no good match anywhere in the control pool. With
a single global $(\min,\max)$, \code{proximity\_minmax}$=1$ is reserved for the single closest
pair observed anywhere in the run, so a unit whose best donor is mediocre still gets a mediocre
score. \textbf{\opt{caliper()} filters on this scale}, which is why it now means the same thing
across specifications with different forest sizes, leaf sizes, or covariate sets --- the raw
scale never did.

\medskip
\noindent\textbf{3. Weight (design weight) --- ``how much does this donor count, within its
own treated unit's block.''} Within the donors actually kept for treated unit $i$, the \emph{raw}
proximities, raised to the power $\gamma$ set by \opt{wpower()}, are renormalised to sum to 1:
\[
\text{weight}(i,j) \;=\; \frac{\text{proximity}(i,j)^{\gamma}}{\sum_{j' \in \mathcal{N}(i)} \text{proximity}(i,j')^{\gamma}},
\qquad \gamma=\text{\opt{wpower()}}\ (\text{default } 1),
\]
where $\mathcal{N}(i)$ is the donor set finally selected for $i$ (after the caliper and/or $k$
cap have been applied). With $\gamma=1$ the weight is proportional to proximity (the only
choice before beta 1.2); with $\gamma=0$ every selected donor gets the same weight (classic
$k$:1 matching); with $\gamma>1$ the weight shifts toward the closest donors and tends to
nearest-neighbour matching as $\gamma$ grows. Raising $\gamma$ lowers the effective sample size,
so check ESS and balance whenever you change it (on the demo panel with $k=50$: mean ESS
$\approx 47$ at $\gamma=0$, $5.6$ at $\gamma=1$, $2.7$ at $\gamma=2$). This is what you pass to \code{[pw=weight]} in the DiD step: a
within-block convex combination that sums to 1 for each treated unit, \emph{not} a 0--1
``how good is this match'' score --- that role belongs to \code{proximity\_minmax}. A treated
unit can have a low \code{proximity\_minmax} for its best donor (poor overlap) while its
\code{weight} column still sums to 1 across whatever donors it did get, so always check
\code{proximity\_minmax} (or the ESS/balance diagnostics) rather than inferring match quality
from \code{weight} alone.

\medskip
\noindent All three are written out by the engine: \code{proximity} and
\code{proximity\_minmax} in both \code{matches.csv} and \code{matched\_panel.csv} (missing, by
construction, on treated rows, which have no proximity to themselves); \code{weight} everywhere
the other two appear, plus as the variable you actually use for estimation.

\subsection{Tuning options (model selection)}
\begin{tabularx}{\textwidth}{@{}lX@{}}
\toprule
Option & Meaning (default)\\
\midrule
\opt{tune}              & turn on cross-validated hyperparameter search --- \textbf{slow} (off)\\
\opt{tuningmetric(str)} & \code{accuracy} $|$ \code{f1\_treated} $|$ \code{combined} (combined)\\
\opt{niter(\#)}         & number of search candidates (40)\\
\opt{cvfolds(\#)}       & cross-validation folds (5)\\
\bottomrule
\end{tabularx}

\noindent\textbf{Tuning never sees the treatment effect.} Hyperparameters are chosen on the
forest's own prediction target (out-of-bag / cross-validated), never on the ATT --- by design.

\subsection{Forest defaults (used when \texttt{tune} is off)}
\begin{tabularx}{\textwidth}{@{}lX@{}}
\toprule
Option & Meaning (default)\\
\midrule
\opt{nestimators(\#)}   & number of trees (1000)\\
\opt{minleaf(\#)}       & minimum samples per leaf (5)\\
\opt{maxfeatures(str)}  & \code{sqrt} $|$ \code{log2} $|$ a fraction like \code{0.5} (sqrt)\\
\opt{maxdepth(str)}     & \code{null} or an integer (null = unlimited)\\
\opt{seed(\#)}          & random seed (42)\\
\bottomrule
\end{tabularx}

\subsection{Output options}
By default the matched panel is loaded into memory and \emph{no files are saved}.

\begin{tabularx}{\textwidth}{@{}lX@{}}
\toprule
Option & Meaning (default)\\
\midrule
\opt{store}        & also save the CSV outputs to \code{./garf\_output}\\
\opt{export(str)}  & also save the CSV outputs to this folder\\
\opt{noload}       & do \emph{not} load the matched panel into memory (produce files only)\\
\bottomrule
\end{tabularx}

\subsection{Plumbing options}
\begin{tabularx}{\textwidth}{@{}lX@{}}
\toprule
Option & Meaning (default)\\
\midrule
\opt{enginedir(str)} & path to the engine folder (else \code{\$GARF\_HOME/engine})\\
\opt{python(str)}    & Python interpreter to use (else the package \code{.venv}, else system Python)\\
\opt{outdir(str)}    & folder for saved files; \emph{implies saving} (\code{./garf\_output})\\
\opt{verbose}        & show the engine's own progress bar / timing\\
\bottomrule
\end{tabularx}

% =====================================================================
\section{What you get back}

\subsection{The on-screen summary}
Treated / control counts, how many treated units were matched, how many distinct donors were
used, mean and maximum absolute standardized mean difference (SMD) before and after matching,
mean ESS per treated unit, Forest Match Quality, and the forest's out-of-bag balanced accuracy and
AUC, with plain out-of-bag accuracy for reference (plus the CV score when tuning is on).

\subsection{Stored results \texttt{r()}}
Everything in the summary is returned, so it is scriptable:
\begin{stata}
garf, treatment(W) outcome(Y) id(id) time(year) t0(2000) tvarying(x1 x2) tinvariant(region)
return list
* r(n_treated) r(n_controls) r(n_treated_matched) r(n_donors)
* r(mean_smd_raw) r(mean_smd_matched) r(max_smd_raw) r(max_smd_matched)
* r(mean_ess) r(fmq_mean) r(fmq_p10) r(fmq_mean_all)
* r(oob_bacc) r(oob_auc) r(oob_score) r(wpower)  [ r(cv_score) if tuned ]
* r(matched_panel)  r(outdir)   // only when files are saved (store/export/outdir)
display "balance improved from " r(mean_smd_raw) " to " r(mean_smd_matched)
\end{stata}

\subsection{Output files (only with \texttt{store} / \texttt{export()} / \texttt{outdir()})}
By default no files are written. When you ask for them:

\begin{tabularx}{\textwidth}{@{}lX@{}}
\toprule
File & Contents\\
\midrule
\code{matched\_panel.csv} & long panel (treated + donor rows) with \code{weight}, \code{proximity\_minmax}, \code{cooccurrence}, \code{fmq}, \code{block\_id}, \code{unique\_id} --- ready for a weighted DiD\\
\code{matches.csv}        & \code{treated\_id}, \code{donor\_id}, \code{rank}, \code{proximity}, \code{proximity\_minmax}, \code{weight}, \code{leaf\_distance}, \code{cooccurrence}\\
\code{blocks.csv}         & one row per treated unit: \code{matched}, \code{n\_donors}, \code{ess}, \code{match\_cooccurrence}, \code{random\_cooccurrence}, \code{best\_cooccurrence}, \code{fmq}, \code{fmq\_twin}, \code{best\_proximity\_minmax}, \code{overlap\_mass}\\
\code{balance.csv}        & standardized mean difference per feature, raw vs matched\\
\code{diagnostics.json}   & config used + all summary numbers + top feature importances\\
\code{features.csv}       & the unit-level design matrix\\
\bottomrule
\end{tabularx}

% =====================================================================
\section{Reading and interpreting the results}\label{sec:interpret}

Every run prints a short summary (and stores the same numbers in \code{r()}). How to read it:

\medskip
\noindent\textbf{Balance --- Mean / Max |SMD| (raw $\to$ matched).} The standardized mean
difference (SMD) is the treated-minus-control gap in a covariate divided by a pooled standard
deviation, so it is comparable across variables. \emph{Raw} is your starting imbalance (treated
vs.\ all controls); \emph{matched} is what GARF achieved (treated vs.\ the proximity-weighted
donors). Matching should pull both toward zero; $|\text{SMD}| < 0.10$ is the usual
``well balanced'' rule of thumb. The \emph{max} matters as much as the mean --- one badly
unbalanced covariate is still a problem.

\medskip
\noindent\textbf{Effective sample size --- Mean ESS per treated.} With proximity weights (and
donor reuse under with-replacement matching), the donors behind a treated unit are not each worth
a full observation. The ESS is their effective count. It is normally \emph{well below} \code{k}
--- that is honest, not a bug --- and you should report it. A very low ESS (say 1--2) means the
effective comparison group is thin: treat those estimates with care, and consider a looser
caliper or a larger \code{k}.

\medskip
\noindent\textbf{Forest Match Quality --- FMQ mean / p10 (all treated).}\label{sec:fmq}
FMQ summarises, on a 0--1 scale, how close each treated unit's donors are \emph{in the
forest's own metric}, relative to what the data make possible. Let $c(i,j)$ be the share of trees in which treated unit $i$ and control $j$
land in the same leaf (column \code{cooccurrence}). For each treated unit,
\[
S_i=\sum_{j\in\mathcal N(i)} \text{weight}(i,j)\,c(i,j),\qquad
\bar c_i=\frac{1}{N_0}\sum_{j\in\text{all controls}} c(i,j),\qquad
S^{\max}_i=\max_{j\in\text{all controls}} c(i,j),
\]
\[
\text{FMQ}_i=\frac{S_i-\bar c_i}{S^{\max}_i-\bar c_i}.
\]
$S_i$ is how often, on average, $i$ shares a leaf with its donors; $\bar c_i$ is what random
donors would achieve; $S^{\max}_i$ is what the closest untreated unit in the data achieves. FMQ is
the share of the \emph{attainable} improvement, from ``a random control'' to ``the best donor
available'', that the matching achieved: $0$ = no closer than random controls, $1$ = every donor as
close as the best one available. Normalising by the best attainable donor keeps the scale
comparable across sample sizes: with tens of thousands of units and small leaves no pair shares a
leaf in most trees, so the earlier version, $(S_i-\bar c_i)/(1-\bar c_i)$, stayed close to 0 even
for good matches. It is still returned as \code{r(fmq\_twin\_mean)} and per block as
\code{fmq\_twin}. If no control is closer than the average one ($S^{\max}_i=\bar c_i$), FMQ is
undefined for that block and counted in \code{r(fmq\_undefined)}. The summary reports the mean
and 10th percentile over \emph{matched} treated units (\code{r(fmq\_mean)}, \code{r(fmq\_p10)})
and the mean over \emph{all} treated units with unmatched ones scored 0 (\code{r(fmq\_mean\_all)});
\code{blocks.csv} and the variable \code{fmq} give the value block by block. It uses only
pre-treatment information, so it can be used freely while designing the match. Read it with care:
\begin{itemize}[leftmargin=1.5em,itemsep=1pt,topsep=2pt]
\item \textbf{It strongly rewards fewer, closer donors.} A smaller \code{k}, a tighter caliper or a
larger \opt{wpower()} systematically raises FMQ, and \code{k(1)} typically gives the highest value. FMQ measures the
\emph{closeness} (bias) side of the design only. \textbf{Never choose a design by FMQ alone}: read
it relative to the number of donors and \emph{together with ESS}, the variance side. A good design
gives up little FMQ for a large gain in ESS. On the demo panel, for example: \code{k(1)} gives
FMQ $0.96$ with ESS $1$; \code{k(5)} $0.75$ with ESS $3.4$; \code{k(50)} $0.63$ with ESS $5.6$;
\code{k(50) wpower(2)} $0.85$ with ESS $2.7$.
\item \textbf{With a caliper}, unmatched treated units drop out of \code{r(fmq\_mean)}, which can
rise simply because the hardest units were excluded. Compare caliper designs on
\code{r(fmq\_mean\_all)}, together with \code{r(n\_treated\_matched)} and ESS.
\item \textbf{Relative, not absolute.} There is no universal ``good'' threshold. Designs that differ
only in \code{k}, caliper, \opt{wpower()} or replacement share the same forest (same covariates,
forest settings and seed) and their FMQ is fully comparable; across covariate sets or forest
settings it is only roughly comparable.
\item \textbf{It sees only what the forest sees} --- closeness in the dimensions that predict
treatment --- so it does not replace the balance table or the pre-treatment outcome path. Blocks
with FMQ close to 0 have donors no closer than random controls: \code{random\_cooccurrence} tells a
\emph{typical} unit (high $\bar c_i$: any control is about as close) from an \emph{isolated} one
(low $\bar c_i$ and low overlap mass).
\end{itemize}
A quick way to see the trade-off:
\begin{stata}
foreach k in 1 5 10 25 50 {
    garf, treatment(W) outcome(Y) id(id) time(year) t0(2000) tvarying(x1 x2) k(`k') noload
    display "k=`k'   FMQ=" %5.3f r(fmq_mean) "   ESS=" %5.1f r(mean_ess)
}
\end{stata}

\medskip
\noindent\textbf{Separability --- Forest OOB balanced accuracy / AUC.} These say how easily the
forest tells treated from control, computed from its \emph{out-of-bag} predictions: each unit is
scored only by the trees that did not see it, so the numbers are honest. The \emph{balanced
accuracy} averages the hit rates of the two groups, so $0.5$ means no better than chance whatever
the share of treated units; the \emph{AUC} is the probability that a random treated unit receives a
higher out-of-bag score than a random control. Values near $0.5$ mean treatment is hard to predict
from the history, which is good for overlap. With rare treatment (a few per cent of units), read
the \textbf{AUC first}: out-of-bag votes then seldom pass the 0.5 threshold, so the balanced
accuracy can stay close to $0.5$ even when the AUC shows the groups are highly separable (on a
panel with 1\% treated units: balanced accuracy $0.57$, AUC $0.98$). The plain out-of-bag \emph{accuracy}
(\code{r(oob\_score)}, in brackets) is kept for reference only: read it against the share of control
units, which a forest reaches by predicting ``control'' for everyone ($0.80$ in the demo panel).
A balanced accuracy or AUC close to 1 (say above $0.95$) is an \textbf{overlap warning}: treated
and controls are almost perfectly distinguishable, the lack-of-positivity problem that afflicts \emph{any} matching method, GARF
included --- it does not mean the tool is broken. When you see it: check that balance still
improves and that ESS is not tiny, and look at the top feature importances in
\code{diagnostics.json} --- if one covariate dominates, ask whether it \emph{mechanically} encodes
treatment or eligibility (a running variable, an id-like field). If so, matching on it is
inherently hard, and you must decide whether it belongs in the covariate set.

\medskip
\noindent\textbf{Pre-trends, not just levels.} Balance on covariate levels is not the whole story
in a panel: a difference-in-differences needs parallel \emph{pre-trends}. Inspect the outcome and
covariate dynamics before $T_0$ --- the matched control path should track the treated path in the
pre-period. This is exactly where matching on the trajectory \emph{shape} is meant to help, and
it is worth a figure in any write-up.

% =====================================================================
\section{The diagnostic protocol: the \texttt{DIAGNOSTICS} folder}\label{sec:diagnostics}

The on-screen summary is a starting point. The package folder \code{DIAGNOSTICS} contains the
full design-stage protocol of the GARF working paper as commented do-files, one per diagnostic.
Every file uses pre-treatment information only, so the whole protocol can be run, and the design
revised, before any outcome is analysed.

\medskip
\noindent\textbf{Run it after your own \code{garf} --- nothing to edit.} \code{garf} stores its whole
design on the matched panel (characteristics \code{\_dta[garf\_*]}: variables, $T_0$, window, $k$,
caliper, \opt{wpower()}, replacement, seed, forest options, the \code{if}/\code{in} sample and the
file it ran on). With the matched panel in memory (or after \code{use garf\_matched.dta}),
\code{cd} to your project folder and run
\begin{stata}
garf, treatment(W) outcome(Y) id(id) time(year) t0(2000) tvarying(x1 x2) store
do "$GARF_HOME/DIAGNOSTICS/00_master.do"
\end{stata}
The diagnostics then describe \emph{your} design. \code{garf} records the \code{.dta} the data were
loaded from, and the diagnostics check that it still holds \code{garf}'s variables and the same
number of rows (sorting or \code{xtset} after loading are fine). If not --- data imported from CSV,
variables created or rows dropped after loading --- save the panel \code{garf} ran on and set its
path first:
\begin{stata}
global GARF_DATA "/path/to/panel.dta"
\end{stata}
\noindent With \opt{store}
(or \opt{export()}/\opt{outdir()}) the diagnostics read \code{garf}'s files directly; without it,
\code{01\_first\_run.do} runs your design once more, with the same seed, to recreate them.
Tables are printed in the log; data and figures go to \code{./garf\_diagnostics}, and at the end
Stata is back in your folder with the matched panel in memory. The placebo backdates your design
by $H=3$ periods, fewer if the data do not reach back far enough for a full window. The
sensitivity grid (19 designs plus yours) is skipped on pools of 20{,}000 units or more unless
\code{global GARF\_GRID on} is set.

\medskip
\noindent\textbf{Run it on the demonstration panel.} Generate the demo panel once
(\code{demo/run\_demo.sh} or \code{.bat}), \code{cd} to a project folder (not the package
folder) and run the same \code{do} command with no matched panel in memory: the manual settings
of \code{00\_settings.do} (the demo panel) are used.

\medskip
\begin{tabularx}{\textwidth}{@{}lX@{}}
\toprule
File & What it does\\
\midrule
\code{00\_settings.do} & Reads the settings from the matched panel in memory; otherwise the manual settings (demo panel)\\
\code{00\_master.do} & Runs every step in order\\
\code{01\_first\_run.do} & Saves \code{raw.dta} (garf's sample) and \code{matched.dta}; finds or recreates \code{garf}'s files\\
\code{02\_balance\_features.do} & $|$SMD$|$ by variable and component (levels vs dynamics), features that worsen\\
\code{03\_variance\_ratios.do} & Variance ratios, raw and matched\\
\code{04\_balance\_over\_time.do} & SMD by period, raw and matched (\code{fig\_balance\_time.pdf})\\
\code{05\_outcome\_path.do} & Pre-treatment outcome path (\code{fig\_outcome\_path.pdf})\\
\code{06\_placebo.do} & Out-of-window placebo: design backdated by $H$ periods (\code{fig\_placebo.pdf})\\
\code{07\_ess\_reuse.do} & ESS per treated unit, donor reuse and concentration of the implied weights\\
\code{08\_overlap.do} & Best-donor normalised proximity and overlap mass (\code{fig\_overlap.pdf})\\
\code{09\_fmq\_blocks.do} & FMQ by block with its components, lowest-FMQ units\\
\code{10\_forest\_importance.do} & Majority-class share and the top features from \code{diagnostics.json}\\
\code{11\_sensitivity\_grid.do} & 19 designs over $k$, caliper, \opt{wpower()}, replacement and seed, plus yours (\code{sens.dta}; slow)\\
\code{12\_frontier.do} & FMQ--ESS frontier (\code{fig\_frontier.pdf})\\
\bottomrule
\end{tabularx}

\medskip
\noindent\textbf{Settings by hand.} To run the protocol without a \code{garf} run first, copy the
folder, edit the manual-settings block of \code{00\_settings.do}, and point Stata at the copy with
\code{global GARF\_DIAG "/path/to/copy"}. To run a single step, run \code{00\_settings.do} and
\code{01\_first\_run.do} once, then the step; the settings stay in force for the session
(\code{macro drop GARF\_*} to start over).

\medskip
\noindent\textbf{Large donor pools.} The widest design of the sensitivity grid is the setting
\code{KWIDE}, default \code{k(1000)} (set \code{global KWIDE} before running to change it). On pools below 1{,}000 untreated
units it is the same as no cap (the forest kernel); on pools of tens of thousands, \opt{nokcap}
can give thousands of donors per treated unit and a matched panel of tens of millions of rows. Each
run of the grid builds the features, fits the forest and computes the proximities from scratch, so
on large data the grid takes hours.

% =====================================================================
\section{The estimation stage: \texttt{garfte}}\label{sec:garfte}

\code{garf} builds the counterfactual: each treated unit $i$ gets a block of donors $j$ with
weights $w_{ij}$ summing to 1. \code{garfte} (new in 1.2) is the step that comes after it. It runs
the package's second Python file, \code{garf\_estimate.py}, with the same Python as \code{garf},
and returns the effect on the treated by event time and overall: a table on screen, an
event-study graph, \code{e(b)} and \code{e(V)}, and files that include the effect of \emph{each}
treated unit in each period. \textbf{The data in memory are never changed.}

\subsection{Syntax and options}
\begin{stata}
* on the matched panel garf left in memory
garfte, method(did|gap) [options]
* on a saved garf output folder
garfte using "folder", method(did|gap) [options]
\end{stata}
\begin{tabularx}{\textwidth}{@{}lX@{}}
\toprule
Option & Meaning\\
\midrule
\opt{method(did|gap)} & required; there is no default, because the two rest on different assumptions\\
\opt{outcome(}\emph{name}\opt{)} & outcome to estimate on; default: the one \code{garf} matched on. Any numeric variable in the matched panel (e.g.\ merged in with \code{merge m:1 id year})\\
\opt{pre(\#)} / \opt{post(\#)} & keep event times $e\ge-\#$ / $e\le\#$; default: all periods in the data\\
\opt{level(\#)} & confidence level (per cent) for the table, the graph's band and the files; default 95, not affected by \code{set level}\\
\opt{nograph} & no event-study graph\\
\opt{outdir(}\emph{str}\opt{)} & folder for the result files; default \code{garf\_output/estimations} (or \emph{folder}\code{/estimations} when \code{garf} saved to a folder, or with \code{using})\\
\opt{enginedir()} / \opt{python()} & as in \code{garf}\\
\bottomrule
\end{tabularx}

\code{garfte} has no \code{if}/\code{in}: dropping rows of a matched panel would change the
counterfactual. On the panel in memory it reads $T_0$, the matching window and the variable names
from the characteristics \code{\_dta[garf\_*]} that \code{garf} stores (they survive
\code{save}/\code{use}); with \code{using} it reads them from the folder's \code{diagnostics.json}.

\subsection{The gap, and the two methods}
Everything starts from the gap between a treated unit and its own GARF counterfactual,
\[
G_{it} \;=\; Y_{it} \;-\; \sum_j \tilde w_{ij}(t)\, Y_{jt},
\]
where $\tilde w_{ij}(t)$ are the \code{garf} weights renormalised over the donors observed in
period $t$ (equal to $w_{ij}$ when no donor is missing).

\medskip
\noindent\textbf{\opt{method(did)} --- matching difference-in-differences.}
\[
\tau_i(e) \;=\; G_{i,T_0+e} \;-\; \bar G_i,\qquad
\bar G_i = \text{mean of } G_{it} \text{ over the matching window } [T_0-L,\,T_0-1]
\]
(or $[T_0-L+1,\,T_0]$ with \opt{includet0}; then $T_0$ is a pre-treatment period and the effects
start at $T_0+1$). Measuring the gap from its window mean removes any level difference left after
matching. Identifying assumption: \emph{parallel trends conditional on the match}. This is the
estimand of Callaway and Sant'Anna's ATT$(g,t)$ with a single cohort, and of trajectory balancing
with demeaning (Hazlett and Xu), with the GARF weights in place of inverse-probability or balancing
weights; the window mean as base averages out period-specific noise, as in synthetic DiD.

\emph{It needs a panel matched without replacement.} With replacement, a donor serving several
treated units makes their blocks dependent. \code{garfte} checks the panel itself: if any donor
appears in more than one block, \opt{method(did)} stops and says so. Rerun \code{garf} with
\opt{noreplacement} and a \opt{k()} the pool can supply (about $k$ controls per treated unit are
needed, otherwise the units matched last get fewer or worse donors). Standard errors are then
clustered by block, which is valid after matching without replacement (Abadie and Spiess, 2022).

\medskip
\noindent\textbf{\opt{method(gap)} --- treated minus its own donors, in levels.}
\[
\tau_i(e) \;=\; G_{i,T_0+e}.
\]
The gap itself, with no differencing, as the post-treatment gap of a synthetic control.
Identifying assumption: \emph{unconfoundedness given the matched pre-treatment history}.
Replacement is allowed: each treated unit keeps its best donors even if another treated unit uses
them too. The gap is credible only if the gaps in the window are small: check the window rows and
\code{e(mean\_window\_gap)}.

\medskip
\noindent\textbf{Report both.} The two estimate the same quantity under different assumptions,
and under fairly general conditions they bracket the true effect (Ding and Li, 2019).

\subsection{Event times, the overall ATT, and what the window rows mean}
With $e = t - T_0$,
\[
\widehat{\text{ATT}}(e) = \frac{1}{N_e}\sum_{i:\,\tau_i(e)\ \text{defined}} \tau_i(e),\qquad
\widehat{\text{ATT}} = \frac{1}{|\mathcal E_{\text{post}}|}\sum_{e\in\mathcal E_{\text{post}}}\widehat{\text{ATT}}(e),
\]
the overall ATT weighting every post-treatment event time shown equally (Callaway and
Sant'Anna's dynamic aggregation). The table labels each period \code{pre} (before the window),
\code{window} (inside it, for the matched outcome), \code{pre (window)} (inside it, for an outcome
that was not matched on) and \code{post}.

\noindent{\color{accent}\textbf{Window estimates are not a pre-trends test.}} With \code{did} they
sum to zero by construction; with \code{gap} the outcome was matched on in the window, so they are
small by construction. Evidence on pre-trends comes from periods \emph{before} the window, from an
outcome that was not matched on, and from the backdated placebo (\code{DIAGNOSTICS/06\_placebo.do}).

\subsection{Standard errors}
Inference is conditional on the matched design (the forest and the weights are taken as given, as
in propensity-score matching). Blocks are the clusters:
\[
\widehat V = \frac{G}{G-1}\sum_{i} d_i d_i',\qquad d_i(e) = \frac{\tau_i(e)-\widehat{\text{ATT}}(e)}{N_e},
\]
with $G$ treated units: the cluster-robust variance of the stacked regression below. When donors are
shared between blocks (\code{gap} on a match with replacement) a donor shock moves several gaps at
once, and \code{garfte} adds, for every shared donor $j$,
\[
\sum_j \Big[A_j(e)A_j(f) - \sum_i a_{ij}(e)a_{ij}(f)\Big]\,\widehat S_j(e,f),\qquad
a_{ij}(e) = \tilde w_{ij}/N_e,\quad A_j(e)=\textstyle\sum_i a_{ij}(e),
\]
the covariance that donor $j$ creates across the blocks that share it; $\widehat S_j$ is the
covariance across periods of donor outcomes around the counterfactual, estimated inside the blocks
that use $j$. Without shared donors the correction is zero. In simulations with heavy donor reuse,
block-clustered intervals without the correction covered the true effect about 90\% of the time at
a nominal 95\%; with the correction, about 95\%. Critical values use Student's $t$ with $G-1$
degrees of freedom; with fewer than about 30 treated units \code{garfte} warns that cluster-type
standard errors can be unreliable.

\subsection{Missing donor-periods, and the checks}
If some donors have no outcome in a period, the weights are renormalised over the donors observed
in that period; the number observed is reported per cell (\code{n\_donors\_obs}) and the number of
renormalised cells is printed. Cells with no donor observed, or with the treated outcome missing,
are skipped; with \code{did}, a treated unit with no gap in the window has no base and is left out.

Before estimating, \code{garfte} checks that every block has one treated unit and that the donor
weights of every block still sum to 1. If not, it stops with \emph{matched panel was modified}:
typically, donors were dropped. Dropping whole periods, or whole blocks (e.g.\ to estimate a
subgroup ATT), is allowed.

\subsection{Files and graph}
Written to \code{garf\_output/estimations/} (or as set above); earlier results of the same method
there are replaced.

\begin{tabularx}{\textwidth}{@{}lX@{}}
\toprule
File & Contents\\
\midrule
\code{garfte\_}\emph{m}\code{\_event.csv} & one row per event time (name, e, year, period, att, se, t, p, lb, ub, n\_treated, mean\_donors\_obs) and a last row \code{ATT}\\
\code{garfte\_}\emph{m}\code{\_units.csv} & the effect of \textbf{each treated unit in each period}: treated\_id, block\_id, year, e, period, tau, pre\_gap\_mean, n\_donors\_obs, n\_donors, ess, fmq\\
\code{garfte\_}\emph{m}\code{\_vcov.csv} & full covariance matrix of the ATT$(e)$ and the ATT\\
\code{garfte\_}\emph{m}\code{\_summary.csv} & key numbers (key,value)\\
\code{garfte\_}\emph{m}\code{\_messages.txt} & the notes and warnings printed on screen\\
\code{garfte\_}\emph{m}\code{.gph} / \code{.pdf} & event-study graph in the style of synthetic-control output: ATT$(e)$ as a line with its \opt{level()}\% confidence band, the matching window shaded, a dashed line at the start of the treatment, the overall ATT and its interval in the subtitle, the legend below the plot. Colours follow the current scheme (\code{set scheme}); in Stata 15+ band and window are transparent (Stata 14: fixed light colours)\\
\bottomrule
\end{tabularx}

\subsection{Stored results}
{\sloppy \code{garfte} is e-class. \code{e(b)} and \code{e(V)} hold ATT$(e)$, named \code{em}$k$ for $e=-k$ and
\code{e}$k$ for $e=k\ge0$, and \code{ATT}, so \code{coefplot} and \code{test} work directly;
\code{e(event)} is the event table as a matrix. Scalars: \code{e(att)}, \code{e(se\_att)},
\code{e(lb\_att)}, \code{e(ub\_att)}, \code{e(p\_att)}, \code{e(level)}, \code{e(df\_r)},
\code{e(N\_treated)}, \code{e(N\_donors)}, \code{e(donors\_reused)}, \code{e(t0)}, \code{e(npre)},
\code{e(window\_first)}, \code{e(window\_last)}, \code{e(mean\_window\_gap)}. Macros: \code{e(cmd)},
\code{e(method)}, \code{e(outcome)}, \code{e(matched\_outcome)}, \code{e(outdir)}.\par}
\begin{stata}
garfte, method(did)
coefplot, vertical keep(e*)
test em10 em9 em8 em7 em6          // periods before the window, jointly zero
\end{stata}

\subsection{Heterogeneity: the units file}
\code{garfte\_}\emph{m}\code{\_units.csv} has one effect per treated unit and period, to relate to
characteristics of the treated units. Two cautions. (i) Unit-level effects are noisy: the donor
average removes noise from the counterfactual, not from the treated unit's own outcome, so their
spread overstates the spread of the true effects; regressing them on characteristics is fine,
reading their variance as heterogeneity is not. (ii) With shared donors they are correlated, and
plain OLS standard errors of such a regression are too small; without shared donors, cluster by
\code{treated\_id} when several periods are used.

\subsection{Replicating \texttt{garfte} with a regression}
The point estimates are those of a stacked regression with block-by-period fixed effects and the
GARF weights (treated rows weight 1, donor weights summing to 1, so every block counts equally):
\begin{stata}
reghdfe Y ibn.year#c.treated_grp [pw=weight], absorb(garf_block#year) vce(cluster garf_block)
\end{stata}
reproduces \opt{method(gap)}; the same regression with the outcome measured from each unit's mean
over the matching window reproduces \opt{method(did)}. With no shared donors the standard errors
agree up to \code{reghdfe}'s small-sample factor. The quick DiD of the next section,
\code{absorb(garf\_unit year)}, uses time effects common to all blocks: close to, but not the same
as, the block-wise comparison of \code{garfte}.

% =====================================================================
\section{Your own estimator: from matched panel to treatment effect}

After a default run the matched (block-expanded) panel is already in memory, and \code{garf} has
added the variables you need for a weighted DiD:
\begin{center}\small
\begin{tabularx}{0.96\textwidth}{@{}lX@{}}
\toprule
Variable & Meaning\\
\midrule
\code{weight} & proximity weight, normalised within each treated block (use as \code{[pw=weight]})\\
\code{proximity} / \code{proximity\_minmax} & raw and 0--1-rescaled match quality (\S\ref{sec:proxcalc}); missing for treated rows\\
\code{treated\_id} / \code{block\_id} & the treated unit each control is matched to\\
\code{treated\_grp} & 1 = treated unit, 0 = selected control\\
\code{control\_sel} & 1 = this row is a selected control (donor)\\
\code{garf\_unit} & numeric id of \code{unique\_id} --- use in \code{absorb()}\\
\code{garf\_block} & numeric id of \code{block\_id} --- use in \code{vce(cluster )}\\
\code{post} & 1 if the period is $\ge T_0$\\
\code{cooccurrence} / \code{fmq} & share of trees in which donor and treated unit share a leaf; block's Forest Match Quality\\
\bottomrule
\end{tabularx}
\end{center}

With the standard treatment coding (1 only in the treated unit's treated years), your treatment
variable is already the treated-times-post indicator, so a proximity-weighted regression of the
outcome on it with unit and time fixed effects returns the ATT directly:
\begin{stata}
preserve                                               // the matched panel replaces your data
garf, treatment(W) outcome(Y) id(id) time(year) t0(2000) tvarying(x1 x2) tinvariant(region)
reghdfe Y W [pw=weight], absorb(garf_unit year) vce(cluster garf_block)
\end{stata}
\code{[pw=weight]} applies the proximity weights; \code{absorb(garf\_unit year)} are the two-way
fixed effects (numeric ids are used because the raw \code{unique\_id}/\code{block\_id} are
strings); clustering on \code{garf\_block} keeps each treated block's rows together.

If instead your treatment variable marks a treated unit in \emph{all} years, build the DiD
interaction from the added dummies:
\begin{stata}
reghdfe Y c.treated_grp#c.post [pw=weight], absorb(garf_unit year) vce(cluster garf_block)
\end{stata}
For horizon-by-horizon (event-study) effects, interact \code{treated\_grp} with event-time
dummies over the post window, still weighted by \code{weight}.

\medskip
\noindent\textbf{Note.} Loading \emph{replaces} the data in memory --- \code{preserve} or save
your working data first, or pass \code{noload} to leave it untouched (add \code{store} to write
the CSVs and read the matched panel later from \code{r(matched\_panel)}).

% =====================================================================
\section{Studying heterogeneous effects}

GARF matches each treated unit to its \emph{own} set of donors and keeps that link explicit:
every control row carries the id of the treated unit it was matched to
(\code{block\_id} / \code{treated\_id}), together with its proximity \code{weight}. That
structure --- a treated unit \emph{plus its personal pool of controls} --- is what makes
heterogeneity analysis natural. You can split the treated population however you like (by size,
sector, region, exposure, \dots) and the correct comparison group follows \emph{automatically}:
each treated unit's matched donors travel with it into its subgroup. A single pooled control
group cannot be sliced this way, because it is not tied to any particular treated unit.

The recipe is two steps:
\begin{enumerate}[leftmargin=1.6em,itemsep=1pt,topsep=2pt]
\item \textbf{Tag each block} by a characteristic of its treated unit, measured in the
      \emph{pre-treatment} period (so the grouping is not itself affected by the treatment).
\item \textbf{Estimate the DiD within each group.} No re-matching is needed --- the controls in a
      group are exactly that group's matched donors.
\end{enumerate}

\noindent\textbf{Conceptual example --- effects by firm size.} Suppose \code{employees} is a
covariate and you want small, medium and large firms separately:
\begin{stata}
* 1. label each block by the treated firm's pre-treatment size
gen s = employees if treated_grp==1 & year < 2016      // treated rows, pre-period only
bysort block_id: egen size_pre = mean(s)               // assign that size to the whole block
gen byte sizeclass = 1 if size_pre <= 25               // small
replace  sizeclass = 2 if size_pre >  25 & size_pre <= 50   // medium
replace  sizeclass = 3 if size_pre >  50 & !missing(size_pre) // big

* 2. within-group ATT: the controls are that group's matched donors
forvalues g = 1/3 {
    reghdfe Y W [pw=weight] if sizeclass==`g', absorb(garf_unit year) vce(cluster garf_block)
}
\end{stata}
Each regression uses only the treated firms in that size band \emph{and their own matched
controls} --- exactly the comparison you want. The same idea extends to a dynamic
(event-study) version by interacting the treated indicator with event-time \emph{within} each
group. One caution: slicing thins the number of treated units per group, so watch the group
counts (\code{tab sizeclass if treated\_grp==1}) --- small groups give noisy estimates and wide
confidence intervals.

% =====================================================================
\section{Worked examples}

\textbf{Default --- load the matched panel into memory, save nothing:}
\begin{stata}
preserve
garf, treatment(W) outcome(Y) id(id) time(year) t0(2000) tvarying(x1 x2) tinvariant(region)
\end{stata}

\textbf{Also keep the CSV files:}
\begin{stata}
garf, treatment(W) outcome(Y) id(id) time(year) t0(2000) tvarying(x1 x2) store
garf, treatment(W) outcome(Y) id(id) time(year) t0(2000) tvarying(x1 x2) export("C:/work/run1")
\end{stata}

\textbf{Files only, leave the data in memory untouched:}
\begin{stata}
garf, treatment(W) outcome(Y) id(id) time(year) t0(2000) tvarying(x1 x2) noload store
\end{stata}

\textbf{Hyperparameter tuning (slow) with a chosen metric:}
\begin{stata}
garf, treatment(W) outcome(Y) id(id) time(year) t0(2000) ///
      tvarying(x1 x2) tinvariant(region) tune tuningmetric(combined) niter(60) cvfolds(5)
\end{stata}

\textbf{Without replacement, as a robustness check:}
\begin{stata}
garf, treatment(W) outcome(Y) id(id) time(year) t0(2000) ///
      tvarying(x1 x2) tinvariant(region) k(10) noreplacement
\end{stata}

\textbf{Custom output folder and an explicit Python interpreter:}
\begin{stata}
garf, treatment(W) outcome(Y) id(id) time(year) t0(2000) tvarying(x1 x2) ///
      outdir("C:/work/garf_run1") python("C:/Python312/python.exe")
\end{stata}

\textbf{Estimation stage on a match without replacement: did, and gap as the bracketing estimate:}
\begin{stata}
garf, treatment(W) outcome(Y) id(id) time(year) t0(2000) tvarying(x1 x2) noreplacement k(3)
garfte, method(did) pre(5) post(5)
garfte, method(gap) pre(5) post(5)
\end{stata}

\textbf{Estimation stage on a saved output folder:}
\begin{stata}
garf, treatment(W) outcome(Y) id(id) time(year) t0(2000) tvarying(x1 x2) noload store
garfte using "garf_output", method(gap)
\end{stata}

\textbf{A subset of the panel (via \texttt{if}):}
\begin{stata}
garf if country=="IT", treatment(W) outcome(Y) id(id) time(year) t0(2000) tvarying(x1 x2)
\end{stata}

% =====================================================================
\section{Hyperparameter tuning}

The forest has hyperparameters (number of trees, minimum leaf size, features per split, depth).
By default GARF uses fixed, sensible values (fast). You can instead let it \emph{choose} them by
cross-validation --- slower, but it can sharpen the metric.

\medskip
\noindent\textbf{How.} Add \code{tune} (and, optionally, the knobs):
\begin{stata}
garf, treatment(W) outcome(Y) id(id) time(year) t0(2000) tvarying(x1 x2) ///
      tune tuningmetric(combined) niter(60) cvfolds(5)
\end{stata}
\code{niter} is how many random hyperparameter combinations are tried; \code{cvfolds} the number
of cross-validation folds. More of either is slower but more thorough.

\medskip
\noindent\textbf{What it optimizes, and why it is safe.} Tuning selects hyperparameters on the
forest's \emph{own} target (how well it predicts \emph{treatment} out of sample), \emph{never} on
the treatment effect. The tool never sees the outcome effect, so there is no way to tune the ATT
--- which is the point, and what keeps the exercise honest.

\medskip
\noindent\textbf{The metric --- \code{tuningmetric()}.} When tuning, models are ranked by one of
\code{accuracy}, \code{f1\_treated}, or \code{combined} (the default). Because treatment is
usually rare, plain \code{accuracy} is a poor criterion: a model that labels everyone a control
can score $0.95$ while being useless for matching. \code{combined} blends, with equal weights, the treated-class F1
with balanced accuracy, rewarding models that actually discriminate treated-like from
control-like units. Prefer \code{f1\_treated} or \code{combined} whenever treated units are a
small share of the sample.

\medskip
\noindent\textbf{When to bother.} Tuning helps most when the default forest is clearly mis-sized
for your data --- very few treated units, very many covariates, or balance that is unstable
across seeds. For a first pass, leave it off; turn it on for the final specification and report
that you did.

% =====================================================================
\section{Changing the engine's defaults}

The Stata options cover the common cases, but the Python engine (\code{engine/garf\_binary.py})
and its \code{config.yaml} expose everything, and you can edit them directly if you run the
engine outside Stata or want a different permanent default.

\medskip
\noindent\textbf{Config values.} \code{engine/config.yaml} lists every setting with a comment
(forest size, caliper, tuning, paths). Editing it changes the defaults for a direct
\code{python garf\_binary.py --config config.yaml} run. Note that the Stata command writes its
\emph{own} config from your options each time, so to change behaviour from Stata you pass options
--- editing \code{config.yaml} does not affect \code{garf}.

\medskip
\noindent\textbf{Forest defaults} (trees, leaf size, features, depth) can be set from Stata with
\code{nestimators()}, \code{minleaf()}, \code{maxfeatures()}, \code{maxdepth()}, or edited in the
config for a direct run. Bigger forests are steadier but slower; smaller \code{minleaf} makes the
proximity sharper (and the overlap warning above more likely).

\medskip
\noindent\textbf{Expert: the model-selection metric.} The \code{combined} score is defined in a
small function near the top of \code{garf\_binary.py}:
\begin{stata}
COMBINED_F1_WEIGHT  = 0.5   # weight on F1 of the treated (positive) class
COMBINED_BAL_WEIGHT = 0.5   # weight on balanced accuracy
def combined_metric(y_true, y_pred): ...
\end{stata}
This is a deliberate plug-in slot. You may want to change it when the default does not suit your
design --- for example, with \emph{very few treated units} you might weight the treated-class F1
more heavily (raise \code{COMBINED\_F1\_WEIGHT}) so the search does not settle for a model that
looks accurate only because it predicts ``control'' almost everywhere; or you might drop in a
different criterion entirely. Only the body of that one function changes; nothing else depends on
it. This is an expert edit --- keep the default unless you have a concrete reason, and report any
change you make.

% =====================================================================
\section{How options map to the engine config}

\code{garf} simply writes a \code{config.yaml} for the engine. The mapping (useful if you also
run the engine directly, or read \code{diagnostics.json}):

\begin{tabularx}{\textwidth}{@{}lX@{}}
\toprule
Stata option & config key\\
\midrule
\opt{treatment} / \opt{outcome} / \opt{id} / \opt{time} & \code{treat\_col} / \code{outcome\_col} / \code{id\_col} / \code{year\_col}\\
\opt{t0} / \opt{npre} / \opt{includet0} & \code{treatment\_year} / \code{n\_pre\_years} / \code{include\_treatment\_year}\\
\opt{tvarying} / \opt{tinvariant} & \code{x\_time\_varying} / \code{x\_time\_invariant}\\
\opt{k} / \opt{noreplacement} & \code{k\_matches} / \code{with\_replacement}\\
\opt{wpower} & \code{weight\_power}\\
\opt{caliper} / \opt{nokcap} & \code{caliper} \quad(\code{nokcap} writes \code{k\_matches: null})\\
\opt{tune} / \opt{tuningmetric} / \opt{niter} / \opt{cvfolds} & \code{tune} / \code{tuning\_metric} / \code{n\_iter\_search} / \code{cv\_folds}\\
\opt{nestimators} / \opt{minleaf} / \opt{maxfeatures} / \opt{maxdepth} / \opt{seed} & \code{n\_estimators} / \code{min\_samples\_leaf} / \code{max\_features} / \code{max\_depth} / \code{random\_state}\\
\opt{outdir} / \opt{verbose} & \code{out\_dir} / \code{progress}\\
\opt{store} / \opt{export} / \opt{noload} & Stata-side only: whether to keep the CSVs, where, and whether to load into memory (not engine keys)\\
\bottomrule
\end{tabularx}

\medskip
\code{garfte} calls \code{garf\_estimate.py} with command-line options: \opt{method()},
\opt{outcome()}, \opt{pre()}, \opt{post()}, \opt{level()} and \opt{outdir()} become
\code{--method}, \code{--outcome}, \code{--pre}, \code{--post}, \code{--level} and \code{--outdir};
\code{using} \emph{folder} becomes \code{--from} \emph{folder}; on the panel in memory, \code{garfte}
exports it to a temporary CSV (\code{--panel}) and passes $T_0$, the window and the variable names
from the characteristics \code{garf} stored. The same estimation runs from a terminal:
\begin{console}
python garf_estimate.py --from garf_output --method did --pre 5 --post 5
\end{console}

% =====================================================================
\section{Troubleshooting}

\begin{tabularx}{\textwidth}{@{}lX@{}}
\toprule
Message / symptom & Fix\\
\midrule
\code{Set the package location first} & \code{global GARF\_HOME "/path/to/GARF-binary"} (or pass \code{enginedir()}).\\
Error about \code{min\_proximity}, or \code{r(oob\_bacc)} missing & An old \code{garf.ado} is running, often a copy in the current folder: \code{which garf}, delete or rename the old copy, then \code{discard}.\\
\code{engine not found} & \code{\$GARF\_HOME}/\code{enginedir()} is wrong; point it at the folder that contains \code{engine/garf\_binary.py}.\\
\code{engine did not produce matched\_panel.csv} & The Python side failed. Run the same call from a terminal to see the error, or pass \code{python("/full/path")} to an interpreter that has \code{pandas}/\code{scikit-learn} (e.g.\ the package \code{.venv}).\\
\code{Columns not found ...} & A variable name doesn't match; check \code{treatment}/\code{outcome}/covariate names against your data.\\
Everything dropped & Check \code{t0}, \code{npre}, and that units have a complete pre-window.\\
Progress bar looks messy in the log & It's off by default in Stata; \code{verbose} turns the engine's own progress back on.\\
Loading replaced my data & By design --- \code{preserve}/save first, or import the panel later from \code{r(matched\_panel)}.\\
Windows path issues & Use forward slashes in \code{GARF\_HOME}/\code{outdir()} (they work on Windows); \code{garf} converts them for the shell call.\\
\bottomrule
\end{tabularx}

\medskip
\noindent\textbf{Messages from \code{garfte}.}\par
\begin{tabularx}{\textwidth}{@{}>{\raggedright\arraybackslash}p{0.34\textwidth}X@{}}
\toprule
Message & Fix\\
\midrule
\emph{no GARF matched panel in memory} & Run \code{garf} (1.2 or newer) first, or \code{garfte using "garf\_output", method(\dots)} on a saved folder.\\
\emph{matched panel was modified} & Donors were dropped from the matched panel. Rerun \code{garf}; dropping whole blocks or whole periods is fine.\\
\emph{method(did) needs a panel matched WITHOUT replacement} & \code{garf \dots, noreplacement k(\#)} with a $k$ the donor pool can supply, or \opt{method(gap)}.\\
\emph{estimation engine not found} & The package is older than 1.2: \code{engine/garf\_estimate.py} is missing.\\
\emph{the estimation engine produced no results} & The Python side failed: run \code{python garf\_estimate.py --from \dots} from a terminal to see the error, or check that the interpreter has \code{scipy}.\\
\bottomrule
\end{tabularx}

\vfill
\begin{center}\small\color{cmt}
GARF-binary · Stata integration · this guide documents \code{garf.ado} and \code{garfte.ado} \textbf{version 1.2}.\\
Federico Fantechi, PhD · Department of Economics, Business and Statistics · University of Palermo · \texttt{federico.fantechi@unipa.it}\\
Comments, bugs and suggestions are very welcome.\\
For the method itself, see the GARF working paper; for the engine, \code{README.txt} in the package.
\end{center}

\end{document}
